\documentclass[11pt]{article}
\usepackage{acl}
\usepackage{times}
\usepackage{latexsym}
\usepackage[T1]{fontenc}
\usepackage[utf8]{inputenc}
\usepackage{microtype}
\usepackage{inconsolata}
\usepackage{graphicx}
\usepackage{url}
\usepackage{xurl}
\usepackage{booktabs}
\usepackage{enumitem}
\usepackage{tabularx}
\usepackage{array}
\usepackage{makecell}
\usepackage{multirow}
\usepackage{float}
\usepackage{amsmath}
\usepackage[section]{placeins}

\setlength\titlebox{5.0cm}
\newcolumntype{Y}{>{\raggedright\arraybackslash}X}
\graphicspath{{./}}

\setlength{\textfloatsep}{6pt plus 1pt minus 2pt}
\setlength{\floatsep}{6pt plus 1pt minus 2pt}
\setlength{\intextsep}{6pt plus 1pt minus 2pt}
\setlength{\dbltextfloatsep}{6pt plus 1pt minus 2pt}
\setlength{\abovecaptionskip}{3pt}
\setlength{\belowcaptionskip}{0pt}
\renewcommand{\arraystretch}{0.96}

\setlist[itemize]{leftmargin=*,itemsep=1pt,topsep=2pt,parsep=0pt,partopsep=0pt}
\setlist[enumerate]{leftmargin=*,itemsep=1pt,topsep=2pt,parsep=0pt,partopsep=0pt}

\title{Socratic-OT: A Grounded Multimodal Socratic Tutor\\
for Anatomy and Neuroscience Learning in Rehabilitation Science}

\author{Vidhyadhari Bandaru \\
  Department of Computer Science \\ and Engineering \\
  University at Buffalo \\
  Buffalo, NY, USA \\
  \texttt{bandaru7@buffalo.edu}
  \And
  Richie M Ilavarapu \\
  Department of Computer Science \\ and Engineering \\
  University at Buffalo \\
  Buffalo, NY, USA \\
  \texttt{richiemo@buffalo.edu}}

\begin{document}
\maketitle

% ─────────────────────────────────────────────────────────────────────────────
\begin{abstract}
We present \textbf{Socratic-OT}, a grounded multimodal tutoring system for
occupational-therapy (OT) students studying anatomy and neuroscience.
Unlike conventional QA assistants that answer directly, Socratic-OT enforces
a \emph{tutor-not-teller} policy: it retrieves supporting evidence exclusively
from a faculty-approved textbook corpus, masks the target answer, and guides
the student through Socratic clues before revealing it.
The system supports both text and image inputs through a five-layer pipeline ---
routing, retrieval, tutoring-policy control, generation, and post-topic
assessment --- orchestrated by a LangGraph finite-state machine.
The text knowledge base is built from all 28 chapters of OpenStax
\emph{Anatomy \& Physiology 2e} (997 retrieval-ready chunks); the multimodal
branch is grounded by a curated set of anatomy diagrams with structured
metadata.
Baseline evaluation on 20 anatomy QA pairs yields RAGAS Faithfulness of
\textbf{0.91} and Answer Relevancy of \textbf{0.97}, confirming strong
retrieval grounding.
A purity audit on 8 scripted tutoring transcripts achieves \textbf{100\%}
Socratic purity --- the masked answer never appears in any clue-phase turn.
A blind multimodal evaluation on uploaded anatomy diagrams achieves
\textbf{66.7\%} structure identification accuracy, with failures concentrated
on cross-sectional and histological views.
We identify retrieval breadth as the primary remaining bottleneck and outline
concrete improvements targeting both retrieval recall and multimodal robustness.
\end{abstract}

% ─────────────────────────────────────────────────────────────────────────────
\section{Introduction}

Large language models have made conversational tutoring widely accessible,
but general-purpose assistants optimize for answering rather than teaching.
For occupational-therapy education, this is a fundamental mismatch.
Anatomy and neuroscience are bottleneck courses for OT programs, characterized
by high content density, strong dependence on 3-D spatial relationships, and
a need for clinical reasoning rather than simple recall.
In these settings, passive answer delivery short-circuits the guided-discovery
process that prepares students for lab and board-style problem solving.

Socratic-OT enforces a system-level \emph{tutor-not-teller} constraint:
the answer is never named in clue-phase turns, retrieval is restricted to
faculty-approved content (OpenStax A\&P~2e), and diagram-based tutoring
applies the same Socratic loop to uploaded anatomy images.
This paper presents the first working prototype and its baseline evaluation
across three dimensions: retrieval grounding (RAGAS), pedagogical purity
(no-leak audit), and multimodal accuracy (blind structure identification).

% ─────────────────────────────────────────────────────────────────────────────
\section{Related Work}

Our design draws on three lines of prior work.

\paragraph{Grounded tutoring.}
SocraticLM \citep{socraticlm} frames tutoring as guided reasoning rather than
answer delivery; we extend this by enforcing the no-reveal constraint
programmatically via a purity guard, rather than relying on soft prompting.

\paragraph{Retrieval-augmented generation.}
We adopt RAGAS \citep{es2024ragas,ragas_docs} for grounding assessment because
its decomposed metrics (faithfulness, relevance, recall, precision) directly
diagnose retrieval and generation failure modes separately.

\paragraph{Multimodal instruction tuning.}
LLaVA \citep{llava,llava16} demonstrates that image-grounded dialogue is
feasible with vision-language models, motivating our VLM cascade design.

Our work differs from Khanmigo \citep{khanmigo} and BoodleBox
\citep{boodlebox,boodlebox_bots}: answer withholding is enforced
programmatically, guaranteeing consistent Socratic behavior regardless of
how the student phrases their request.

% ─────────────────────────────────────────────────────────────────────────────
\section{Data}

\subsection{Text Knowledge Base}

The text corpus is derived from OpenStax \emph{Anatomy \& Physiology 2e}
\citep{openstax_ap2e}, the standard approved reference for OT anatomy courses.
Our pipeline: (1)~extracts full text from all 28 chapters via PDF parsing;
(2)~detects section boundaries and generates section-level files;
(3)~cleans page markers, duplicated fragments, and formatting artifacts;
(4)~assigns topic labels (\emph{muscle}, \emph{nervous system}, etc.) to each
section;
(5)~chunks each cleaned section into overlapping paragraphs of approximately
341 words with 40-word overlap;
(6)~stores each chunk with source identifier, chapter, section title, topic
label, and keyword summary.

This yields \textbf{997 retrieval-ready chunks}, embedded with
\texttt{all-MiniLM-L6-v2} \citep{allminilm_model} and indexed in
ChromaDB \citep{chroma_docs}.
Table~\ref{tab:kb-assets} summarizes the full knowledge-base inventory.
Figure~\ref{fig:chunk-sample} shows a representative chunk to illustrate
the granularity and structure of the retrieved evidence.

\begin{figure}[t]
\centering
\fbox{%
\begin{minipage}{0.93\columnwidth}
\footnotesize
\textbf{Chunk ID:} CH0290 \quad
\textbf{Chapter:} 10 \quad
\textbf{Topic:} \texttt{muscle}\\[2pt]
\textbf{Section:} 10.1 Overview of Muscle Tissues\\[4pt]
The body contains three types of muscle tissue: skeletal, smooth, and
cardiac. Contraction begins when actin is pulled by myosin after binding
sites are exposed via Ca\textsuperscript{2+} interaction with troponin
and tropomyosin\ldots
\end{minipage}}
\caption{Representative text chunk from the knowledge base (1 of 997).
Chunks are stored with source identifiers, chapter, section, topic label,
and keyword summary.}
\label{fig:chunk-sample}
\end{figure}

\begin{table}[t]
\centering\small
\begin{tabularx}{\columnwidth}{@{}l r@{}}
\toprule
\textbf{Asset} & \textbf{Count} \\
\midrule
Source chapters & 28 \\
Retrieval-ready text chunks & 997 \\
Embedding model & \texttt{all-MiniLM-L6-v2} \\
Vector store & ChromaDB \\
\bottomrule
\end{tabularx}
\caption{Knowledge-base assets at Milestone~2.}
\label{tab:kb-assets}
\end{table}

\subsection{Image Knowledge Base}

The visual knowledge base consists of a curated collection of anatomy diagrams
aligned with the OT curriculum scope, spanning major anatomical systems:
nervous tissue, brain organization, peripheral nerve networks, and skeletal
muscle structure.
Each diagram is paired with a structured metadata record containing: image ID,
canonical structure name, known aliases, topic label, body region, functional
description, and a related text topic.
Table~\ref{tab:image-metadata} shows an example record.

\begin{table}[t]
\centering\footnotesize
\begin{tabularx}{\columnwidth}{@{}lY@{}}
\toprule
\textbf{Field} & \textbf{Value} \\
\midrule
Image ID & IMG001 \\
Structure name & Neuron \\
Aliases & nerve cell; neuron structure \\
Topic & nervous tissue \\
Body region & nervous system \\
Function & basic neural signaling structure \\
\bottomrule
\end{tabularx}
\caption{Example image metadata record (IMG001).}
\label{tab:image-metadata}
\end{table}

This metadata \emph{bridges} uploaded diagrams to the text KB, grounding
image-path tutoring in the same retrieval pipeline as text-path tutoring.

\subsection{Evaluation Data and Transcripts}
\label{sec:transcripts}

For purity auditing and qualitative analysis we construct eight scripted
evaluation transcripts covering the full scenario space of the system:
five text-path (normal flow, early correct guess, repeated DONT\_KNOW,
out-of-KB rejection, wrong-named comparison clue) and three image-path
(in-KB diagram, non-anatomy domain guard, unlabeled 3-D render with fuzzy
KB match).

\medskip
\noindent\textbf{Text-path excerpt} (eval\_s2: student guesses correctly after
the second clue):

\begin{small}
\begin{tabbing}
\hspace{1.5em}\=\kill
\textbf{S:}\> What lobe of the cerebral cortex is primarily \\
            \> responsible for visual processing? \\[3pt]
\textbf{T:}\> Great question, let's explore this together! This lobe \\
            \> is located on the \emph{posterior} aspect of the cerebrum, \\
            \> near the midline, above the tentorium cerebelli. \\
            \> Can you name the lobe? \\[3pt]
\textbf{S:}\> Is it the frontal lobe? \\[3pt]
\textbf{T:}\> Not quite --- that's a different lobe. This lobe processes \\
            \> and integrates sensory information. What is it called? \\[3pt]
\textbf{S:}\> The occipital lobe? \\[3pt]
\textbf{T:}\> Yes, exactly --- \textbf{Occipital lobe}! The occipital lobe \\
            \> is the region primarily responsible for processing visual \\
            \> information and interpreting visual stimuli\ldots
\end{tabbing}
\end{small}

\medskip
\noindent\textbf{Text-path excerpt} (eval\_s3: student exhausts clue budget
and requests the answer):

\begin{small}
\begin{tabbing}
\hspace{1.5em}\=\kill
\textbf{S:}\> What nerve network originates from the cervical and upper \\
            \> thoracic spinal roots and supplies the upper limb? \\[3pt]
\textbf{T:}\> Great question, let's explore this together! This network \\
            \> arises from the cervical and upper thoracic spinal regions \\
            \> and is responsible for supplying the upper limb with a \\
            \> wide range of motor and sensory functions. What is it? \\[3pt]
\textbf{S:}\> I don't know. \\[3pt]
\textbf{T:}\> No worries, let me give you another angle. This network \\
            \> enables movements from simple grasping to fine motor \\
            \> control. What is it called? \\[3pt]
\textbf{S:}\> I have no idea, I'm stuck. \\[3pt]
\textbf{T:}\> That's okay --- let's try a different clue. This network \\
            \> is more localized than the sympathetic system --- focused \\
            \> on the cervical and upper thoracic region. What is it? \\[3pt]
\textbf{S:}\> I really don't know, I give up. \\[3pt]
\textbf{T:}\> No worries! The answer is \textbf{brachial plexus} --- a \\
            \> network of nerves from the neck and upper back controlling \\
            \> the shoulder, elbow, wrist, and hand\ldots
\end{tabbing}
\end{small}

\medskip
\noindent\textbf{Image-path excerpt} (eval\_img2: non-anatomy image rejected
before tutoring begins):

\begin{center}
  \includegraphics[width=0.55\columnwidth]{eval_img2_upload.png}
\end{center}

\begin{small}
\begin{tabbing}
\hspace{1.5em}\=\kill
\textbf{S:}\> \textit{[uploads image above]} What is this diagram showing? \\[3pt]
\textbf{T:}\> Hmm, that doesn't look like an anatomy diagram to me! \\
            \> I'm designed to help with anatomy and neuroscience images \\
            \> --- things like muscles, nerves, or brain diagrams. \\
            \> Try uploading an anatomy diagram and I'll be happy \\
            \> to work through it with you.
\end{tabbing}
\end{small}

\medskip
\noindent\textbf{Image-path excerpt} (eval\_img3: unlabeled 3-D neuron render
matched to KB via fuzzy alias):

\begin{center}
  \includegraphics[width=0.55\columnwidth]{eval_img3_upload.png}
\end{center}

\begin{small}
\begin{tabbing}
\hspace{1.5em}\=\kill
\textbf{S:}\> \textit{[uploads image above]} I see a cell with branches --- \\
            \> what structure is this? \\[3pt]
\textbf{T:}\> Oh nice, let's dig into this diagram together! This is a \\
            \> peripheral branching network of structures. Can you tell \\
            \> me which cell this is? \\[3pt]
\textbf{S:}\> Is it a Schwann cell? \\[3pt]
\textbf{T:}\> Not quite --- that's a different cell. The branches here \\
            \> provide support and insulation for a basic neural signaling \\
            \> structure found in the PNS. What cell is this? \\[3pt]
\textbf{S:}\> Maybe a neuron? \\[3pt]
\textbf{T:}\> Yes, exactly --- \textbf{neuron}! A neuron is a specialized \\
            \> cell that plays a crucial role in transmitting information \\
            \> throughout the nervous system\ldots
\end{tabbing}
\end{small}

% ─────────────────────────────────────────────────────────────────────────────
\section{Solution Architecture}

Socratic-OT is organized into five runtime layers (Figure~\ref{fig:architecture}).
Text and image inputs initialize the hidden target differently, but both
converge to the same retrieval-grounded tutoring controller.

\begin{figure}[t]
  \centering
  \includegraphics[width=\columnwidth]{architecture_diagram-3.png}
  \caption{Socratic-OT runtime architecture. Text and image inputs share the
  same retrieval-grounded tutoring controller after target initialization.}
  \label{fig:architecture}
\end{figure}

\subsection{Routing and Retrieval}

Every student input passes through a routing layer that classifies it as
text, image, or follow-up.
For text inputs, the retriever computes an \texttt{all-MiniLM-L6-v2} embedding
and retrieves the top-5 chunks from ChromaDB by cosine similarity.
A metadata-aware reranking step then prioritizes chunks whose topic label
aligns with the inferred query intent; adjacent chunks from the same section
may be merged before final prompting to improve context coherence.

For image inputs, the VLM module runs a cascade:
Groq \texttt{llama-4-scout-17b} (primary) $\rightarrow$ LLaVA-NeXT on GPU
(when available) $\rightarrow$ GPT-4o (fallback) \citep{gpt4o,llava16}.
The VLM returns an IS\_ANATOMY flag, predicted structure name, topic, and
confidence.
If IS\_ANATOMY is false, a polite domain-guard rejection is issued immediately
and no tutoring begins.
Otherwise, the predicted structure is matched against KB metadata (exact name
or alias), and the matched record provides both the canonical target and the
text-retrieval topic for subsequent tutoring.

\subsection{Tutoring Policy Controller}

The tutoring policy layer is implemented as a LangGraph finite-state machine
\citep{langgraph_docs} with phases:
\textsc{rapport} $\to$ \textsc{clue}(1--3) $\to$ \textsc{reveal} $\to$
\textsc{post\_reveal\_wait} $\to$ \textsc{quiz} $\to$ \textsc{done}.
It enforces four system-level invariants:

\paragraph{Answer masking.}
At session start, the system constructs a forbidden-term set from the masked
answer plus LLM-generated aliases (e.g., ``neuron'' $\to$
\{``neuron'', ``nerve cell'', ``neural cell''\}).
A purity guard scans every generated response and triggers regeneration if any
forbidden term appears.
To allow comparison clues --- where the tutor must name the student's wrong
guess to contrast it with the hidden target --- the student's own attempt
tokens are excluded from the leak check.

\paragraph{Attempt classification.}
An LLM judge classifies each student response as \textsc{correct},
\textsc{partial}, \textsc{wrong\_named}, \textsc{dont\_know}, or
\textsc{other}.
For image-path tutoring, a deterministic fast-path classifies known
biologically-proximate wrong guesses (e.g., ``astrocyte'' or ``glial cell''
when the answer is ``neuron'') using a pre-specified list in the image
metadata, bypassing the LLM to prevent classifier ambiguity on
anatomically similar structures.

\paragraph{Clue dimension selection.}
Clue content is shaped by the attempt classification.
A \textsc{wrong\_named} response triggers a \emph{comparison} clue that
explicitly names the student's wrong guess and explains how it differs from
the hidden target --- without naming the target.
A \textsc{dont\_know} response triggers a progressively stronger hint without
a comparison anchor.

\paragraph{Domain guard.}
An LLM-based scope classifier determines whether the query concerns human
anatomy or physiology.
Out-of-domain questions receive a polite rejection rather than a hallucinated
in-domain answer.

\subsection{Generation and Assessment}

The generation layer uses Groq-hosted Llama~3.1~8B Instruct as the primary
text backend, with GPT-4o-mini as fallback \citep{groq_docs,gpt4omini_docs}.
After each tutoring topic, a lightweight post-topic assessment module presents
the student with three mastery-check questions drawn from the retrieved context.
Session state tracks weak topics and wrong guesses for within-session revisit.

\subsection{Design Rationale}

The retrieval layer is domain-agnostic: swapping the ChromaDB collection
ports the tutor to a new subject without architecture changes.
The policy layer enforces pedagogical constraints independently of the
generation layer --- the LLM generates clue text, but \emph{whether} to
reveal the answer is decided entirely by the state machine, making
Socratic behavior robust to unusual student phrasing or prompt injection.

% ─────────────────────────────────────────────────────────────────────────────
\section{Experiments and Results}

\subsection{Groundedness: RAGAS Evaluation}
\label{sec:ragas}

We evaluate retrieval-augmented generation quality using RAGAS
\citep{es2024ragas} on 20 anatomy QA pairs spanning all 28 KB chapters.
For each question, the system retrieves the top-5 context chunks and generates
an answer conditioned strictly on that retrieved context.
We report four metrics: Faithfulness, Answer Relevancy, Context Recall, and
Context Precision.

\begin{table}[t]
\centering\small
\begin{tabularx}{\columnwidth}{@{}l c c c@{}}
\toprule
\textbf{Metric} & \textbf{Score} & \textbf{Target} & \textbf{Status} \\
\midrule
Faithfulness      & 0.9135 & $\geq$0.90 & Pass \\
Answer Relevancy  & 0.9694 & $\geq$0.85 & Pass \\
Context Recall    & 0.6471 & $\geq$0.80 & Below target \\
Context Precision & 0.7792 & $\geq$0.80 & Slightly below \\
\bottomrule
\end{tabularx}
\caption{Baseline RAGAS results on 20 anatomy QA pairs.}
\label{tab:ragas-results}
\end{table}

\textbf{Faithfulness (0.91)} and \textbf{Answer Relevancy (0.97)} confirm
the system stays grounded and on-topic --- it does not hallucinate anatomy
facts and directly addresses each question.
\textbf{Context Recall (0.65)} is the weakest result: broad questions
spanning multiple chapters (e.g., \emph{``Describe the dorsal
column--medial lemniscal pathway''}) require evidence distributed across
several section chunks, but a single semantic query surfaces only the most
proximate unit.
\textbf{Context Precision (0.78)} falls slightly below target due to
off-topic chunks retrieved via lexical overlap with unrelated sections ---
both metrics point to retrieval breadth, not generation quality, as the
bottleneck.

\subsection{Socratic Purity Audit}

We define \emph{Socratic purity} as the property that the masked answer does
not appear in any \textsc{clue}-phase tutor turn.
We audit this by scanning all 8 evaluation transcripts and checking
case-insensitively whether the masked answer string appears in any tutor turn
with phase~$=$~\textsc{clue}.

\textbf{Result: 8/8 = 100\% purity.}
This holds across adversarial cases including biologically-proximate wrong
guesses (e.g., ``astrocyte'' when the answer is ``neuron''), where the
purity guard must allow the student's guess to be named in a comparison clue
while still suppressing the hidden target.

\paragraph{Clue relevance (qualitative).}
Purity verifies the answer is not leaked; it does not verify the clue is
useful.
We therefore manually assessed all clue-phase tutor turns across the 8
transcripts on a 1--5 rubric (1~=~generic or off-topic;
3~=~partially relevant; 5~=~grounded, specific, and clearly narrowing
toward the answer without revealing it).
Clues for focused single-topic questions consistently rated 4--5, reflecting
strong retrieval grounding.
Clues for broad multi-chapter questions rated 3--4, directly tracking the
Context Recall gap (Section~\ref{sec:ragas}) --- when the retriever misses
relevant chunks, the clue loses specificity.
This qualitative pattern motivates the hybrid retrieval improvement planned
for future work; a formal student-cohort annotation is reserved for
final-milestone evaluation.

\subsection{Multimodal Blind Test}

We evaluate the multimodal branch on a blind evaluation set of uploaded
anatomy diagrams, each passed to the VLM without labels or KB hints.
The VLM must identify the anatomical structure from visual content alone;
a prediction is scored correct if the expected structure name or any
registered alias appears in the predicted string.
Table~\ref{tab:vlm-results} shows per-image results.

\begin{table}[t]
\centering\small
\begin{tabularx}{\columnwidth}{@{}XXc@{}}
\toprule
\textbf{Image} & \textbf{Predicted} & \textbf{Correct} \\
\midrule
Neuron structure          & neuron                  & \checkmark \\
Brain lobes               & brain lobes             & \checkmark \\
Spinal cord cross-section & skeletal muscle fiber   & $\times$   \\
Nervous system overview   & nervous system overview & \checkmark \\
Brachial plexus           & brachial plexus         & \checkmark \\
Skeletal muscle fiber     & brachial plexus         & $\times$   \\
\midrule
\multicolumn{2}{@{}l}{\textbf{Overall}} & \textbf{4/6 = 66.7\%} \\
\bottomrule
\end{tabularx}
\caption{Blind-test VLM structure identification (backend: Groq \texttt{llama-4-scout-17b}).}
\label{tab:vlm-results}
\end{table}

\textbf{Result: 4/6 = 66.7\%} (target $\geq$80\%).
Both failures involve histological and cross-sectional views.
The spinal cord cross-section and skeletal muscle fiber microscopy images share
visual features --- elongated fiber structures, staining patterns --- that the
vision model conflates.
Diagram-style and overview images, which have clearer spatial landmarks, are
all identified correctly.

% ─────────────────────────────────────────────────────────────────────────────
\section{Discussion}

The results confirm that \textbf{generation quality is strong and retrieval
breadth is the primary bottleneck.}
The 100\% purity score validates separating pedagogical policy (the state
machine) from generation (the LLM): the answer-masking invariant holds even
under adversarial conditions.

\paragraph{Retrieval limitations and planned fix.}
Context Recall (0.65) exposes the limitation of top-$k$ semantic retrieval
for multi-chapter questions: a single embedding query surfaces only the
most lexically proximate chunk, leaving evidence distributed across other
sections unretrieved.
Context Precision (0.78) reflects a secondary issue: a small proportion of
retrieved chunks are off-topic due to lexical overlap with unrelated sections.
Our planned fix is a hybrid retrieval pipeline combining BM25 sparse retrieval
and dense embedding search, merged by reciprocal rank fusion, with a
topic-label filter applied post-merge.
This is expected to lift Context Recall toward the $\geq$0.80 target without
sacrificing the Faithfulness already achieved.

\paragraph{Multimodal limitations and planned fix.}
The VLM accuracy of 66.7\% is interpretable: both failures share a specific
visual failure mode --- confusion of elongated fiber textures in
cross-sectional and histological views --- while diagram-style and overview
images with clear spatial landmarks are all correctly identified.
The fix is two-pronged: (1)~enrich image metadata with histological cues
(staining pattern, cross-section orientation, distinctive landmarks) to
compensate when VLM confidence is low; (2)~introduce confidence-based
routing so that low-confidence predictions are escalated to GPT-4o
\citep{gpt4o} rather than accepted from the primary model.

\paragraph{Scope of current evaluation.}
Several components are implemented but not yet experimentally validated.
Session memory and weak-topic revisit operate within a single session;
cross-session persistence, mastery assessment effectiveness, and domain
generalizability (KB swap to a new subject) are reserved for
final-milestone evaluation and are not claimed as demonstrated results here.

% ─────────────────────────────────────────────────────────────────────────────
\section{Conclusion}

Socratic-OT demonstrates that a pedagogically controlled, retrieval-grounded
multimodal tutoring system is feasible using open-source components and
free-tier APIs.
Socratic purity holds at 100\% across all evaluated scenarios including
adversarial cases, and generation quality meets targets (Faithfulness 0.91,
Answer Relevancy 0.97).
The two primary gaps --- retrieval breadth (Context Recall 0.65) and
multimodal robustness on histological images (66.7\%) --- are
well-characterized and have concrete improvement paths: hybrid BM25 + dense
retrieval and confidence-based GPT-4o routing respectively.
Future work will validate these improvements and evaluate cross-session
memory, mastery assessment, and domain transfer at final milestone.

% ─────────────────────────────────────────────────────────────────────────────
\bibliography{custom}

\end{document}
